% !TEX root = ../nc-quadric-moduli.tex
% 2026-11-22 / v1.9（相対定義とflatness剛性補題のB移管）
\section{内在的な相対対象と基底変更}
\label{sec:moduli-relative}

本節では $k$ を標数零の代数閉体とし、まず $k$ 上局所有限型の
スキームの big fppf site を用いる。この範囲は非被約スキームと
有限型の局所 Artin 環を含む。任意の noetherian $k$-代数に対する
連続拡張や極限保存性を、本節の構成だけから結論しない。
とくに finite type の基底での局所 noetherian 性と、任意の
noetherian 基底変更での strong noetherian 性を区別する。

\subsection{Abel 圏の基底変更を指定する}

局所有限型の affine スキーム $U=\operatorname{Spec}A$ 上で、
$A$-線形 noetherian Abel 圏を $\mathcal C_U$ とし、
$\mathcal A_U=\operatorname{Ind}\mathcal C_U$ と置く。
$A\to A'$ に対し $(\mathcal A_U)_{A'}$ は
$\mathcal A_U$ 内の $A'$-加群対象の圏、すなわち
$A'$ の作用を備えた対象とその作用を保つ射の圏である。
この構成において引戻しは
$A'\otimes_A-$ であり、一般には右完全、$A'$ が $A$ 上平坦なら完全である。
本稿でいう圏の基底変更は、Hom 加群を単に tensor して対象集合を固定する
線形圏の基底変更ではない。
この区別と基本的な随伴は
\cite[\S4]{LowenVanDenBergh2006}、任意の有限型 affine 基底に必要な
加群対象の定式化と flat descent は
\cite[\S7.1, Definition 7.3, Remarks 7.4--7.5]{AHLHExistence}
を参照する。後者は Artin--Zhang の abstract Hilbert schemes の
B2.2, B3.16, B6.3, C8.6 を明記して用いている。

\begin{definition}[基底変更を備える Abel 圏の族]
\label{def:relative-abelian-system}
$T$ 上の Abel 圏の族とは、$T$ 上の locally finite type affine
$U$ に $\mathcal C_U$ を割り当てる pseudofunctor であって、次を満たすものとする。
\begin{enumerate}
\item 各 $\mathcal C_U$ は $A$-線形で本質的に小さい noetherian Abel 圏であり、
  その Ind 圏は局所 noetherian Grothendieck 圏である。
\item affine の間の射 $U'=\operatorname{Spec}A'\to U$ について
  \[
    \mathcal A_{U'}\simeq(\mathcal A_U)_{A'}
  \]
  という基底変更同値があり、引戻しは $A'\otimes_A-$ に対応する。
  これらの比較同値は合成と恒等射に対して coherent である。
\item 対象、射、完全列について fppf descent が成立する。
  とくに $\mathcal C$ は対象と全ての射を保持する stack である。
\item $\operatorname{Hom}_{\mathcal C_U}(E,F)$ は有限 $A$-加群であり、
  affine flat base change と可換する。
\end{enumerate}
射はこの系の間の、基底変更と可換する線形完全同値とする。
\end{definition}

この定義は圏の局所表示を指定するものではない。
基底変更同値は、族という語の意味を固定するためのデータであり、
射はその coherent な同型を許す。
有限型 affine 基底変更での局所 noetherian 性は
Artin--Zhang B6.3 の範囲内である。
任意の noetherian 基底変更へ、この結論を無条件に拡張しない。

$\mathcal C_U$ の対象 $E$ が $A$ 上平坦とは、
$-\otimes_AE:\operatorname{Mod}A\to\mathcal A_U$ が完全であることをいう。
本稿では
\[
 \mathcal D_U=D(\mathcal A_U)^c
\]
を標準の導来 enhancement における compact 対象の圏とする。
この圏は $\mathcal C$ から決まり、別の dg enhancement を標識に加えない。
以下の導来条件は $T$ の fppf 局所で課す。
すなわち fppf 被覆 $T_i\to T$ が存在して、その各 $T_i$ への
全ての本節の affine 基底変更 $U$ 上で、$D(\mathcal A_U)$ は
compactly generated であり、$\mathcal D_U$ は $D^b(\mathcal C_U)$ に含まれ、
相対 smooth proper な $\operatorname{Perf}A$-線形圏となり、導来基底変更
\[
 \mathcal D_U\otimes_{\operatorname{Perf}A}\operatorname{Perf}A'
       \simeq\mathcal D_{U'}
 \tag{RP-BC}\label{eq:relative-perfect-basechange}
\]
と整合することを要求する。
$\mathcal O$ はこれらの局所 chart 上でこの compact 部分に属するものとする。
chart の選択は標識に含めない。
本節の stack 性の証明には、貼合せ後の大域的な
$D(\operatorname{Ind}\mathcal C_U)$ が局所導来圏の descent と一致するという
追加主張を用いない。定義そのものが fppf 局所条件であることを用いる。
特異な基底上で $\mathcal D_U=D^b(\mathcal C_U)$ としたり、
compact 部分だけに bounded heart を課したりしない。
また全ての底上平坦 coherent 対象が compact であることを
今回の定義には追加しない。この一致は別の相対認識問題である。
体上では、以下の foundation による議論で
$\mathcal D_k=D^b(\mathcal C_k)$ が従う。

以下の Hom、双対、導来 tensor はこの標準 enhancement で解釈する。
(\ref{eq:relative-perfect-basechange}) は導来圏の基底変更であり、
heart の通常の引戻しが全ての基底変更で完全であるという主張ではない。

\begin{definition}[相対 Serre 条件]
\label{def:relative-serre-condition}
上の族が相対 Serre 条件を満たすとは、前記の局所 chart の各 affine $U$ 上で
次が成立し、任意の本節の基底変更と整合することをいう。
\begin{enumerate}
\item compact 対象 $E,F$ の間の $\mathbf R\operatorname{Hom}(E,F)$
  は $A$ 上 perfect で、導来基底変更と可換する。
\item $\mathcal D_U$ に自己同値 $S_U$ が存在し、
  functorial な同型
  \[
   \mathbf R\operatorname{Hom}(E,F)^\vee_A
       \simeq\mathbf R\operatorname{Hom}(F,S_UE)
  \]
  を与える。
\end{enumerate}
Serre 関手および duality の選択自体は標識に含めない。
二つの選択の比較は、表現する双対 functor の Yoneda 一意性による。
\end{definition}

\subsection{相対的な強い\texorpdfstring{$\mathcal H_2$}{H2}条件}

相対 quadratic AS4 被覆とは、局所 affine 基底 $A$ 上の次数付き環
$\mathcal B=T_A(V)/(I)$ であって、$\operatorname{rk}V=4$、
$I\subset V^{\otimes2}$ が階数六の部分束、
各次数が有限局所自由、左右 noetherian、
全幾何 fiber が型 $(1,4,6,4,1)$ の AS 正則環となるものをいう。
加えて、標準の左右 Koszul 複体は全次数で完全であり、
基底変更後にも同じ完全列を与えるものとする。
ここで幾何 fiber の定義には、記述された有限個の係数を含む体からの
体拡大で作る次数付き環を用いる。
被覆の二次の正則正規元 $g$ は、左右乗法の単射性と正規性に加え、
$\mathcal B/(g)$ の次数別局所自由性を要求する。
従って切断の完全列は任意の本節の基底変更と両立する。
固定した相対切断環への被覆の有限判定は、分類主論文 \cite[相対AS被覆の有限判定]{QuadricIntrinsic2026} を参照する。
平方根を選ばず、その存在だけを要求する部分関手の代数性は別問題である。

\begin{definition}[有限型基底上の内在的 $\mathfrak H_2$]
\label{def:relative-h2}
$\mathfrak H_2(T)$ の対象は、定義~\ref{def:relative-abelian-system} の族
$\mathcal C$ と対象 $\mathcal O\in\mathcal C_T$ の対で、以下を満たすものとする。
\begin{enumerate}[label=(RQ\arabic*)]
\item $\mathcal O$ は基底上平坦で、fppf 局所 chart の compact 導来部分は
  (\ref{eq:relative-perfect-basechange}) と相対 Serre 条件を満たす。
  これらの chart の全ての affine 基底変更上で
  $\mathbf R\operatorname{Hom}(\mathcal O,\mathcal O)\simeq A[0]$。
\item 各幾何 fiber の対 $(\mathcal C_{\bar t},\mathcal O_{\bar t})$ は、
  分類主論文 \cite{QuadricIntrinsic2026} で定める体上の $\mathcal H_2$ に属する。
  特に標識は同論文の Gabriel--Krull 商の意味で rank one torsion-free であり、
  Serre 条件は fiber の全 $D^b(\mathcal C_{\bar t})$ に成立する。
  fiber 圏は module-object 基底変更、または下記の全軌道環の
  $\qgr$ によって計算し、その一致を要求する。
\item $T$ の fppf 局所で、系全体の線形完全自己同値 $s$ が存在する。
  前記の chart と共通細分を取った局所 compact 導来部分で
  $s^2\simeq S^{-1}[2]$ であり、全ての affine 基底変更上で
  $(\mathcal O,s)$ は Artin--Zhang ample である。
  ここで生成条件は有限直和からの全射、Hom の条件は各全射に対する
  十分大きなシフトでの全射性として課す。
\item 同じ局所的な $s$ について、前記の相対 AS4 環 $\mathcal B$ と
  二次の正則正規元 $g$ が存在し、全次数で
  \[
   \bigoplus_{n\ge0}\operatorname{Hom}(\mathcal O,s^n\mathcal O)
     \simeq\mathcal B/(g)
  \]
  となる。この同型は任意の本節の基底変更と可換する。
\end{enumerate}
射は pointed exact 同値 $(F,u)$、二射は
$u'\eta_{\mathcal O}=u$ を満たす自然同型とする。
$s$、根同型、被覆、切断元、軌道環の表示のいずれも対象の標識ではない。
\end{definition}

この定義は「全幾何 fiber が $\mathcal H_2$ に属する」ことだけを課す定義より強い。
fppf 局所に存在すべき平方根と AS 被覆は、非被約底の変形方向を含めて要求される。
体上の分類と fiberwise な判定だけで、この相対条件を確認したとは扱わない。

\begin{lemma}[Abel flatness と標識付き二射の剛性]
\label{lem:relative-h2-flat-center}
定義~\ref{def:relative-h2} の局所族について、次が成り立つ。
\begin{enumerate}
\item $\mathcal C_U$ は Lowen--Van den Bergh の意味で $A$ 上 flat な Abel 圏である。
\item $\operatorname{Aut}(\operatorname{Id}_{\mathcal C_U})=A^\times$。
  特に pointed 一射の二自己同型は自明である。
\end{enumerate}
\end{lemma}
\begin{proof}
$A$-線形完全自己同値は tensor と可換するから、
$s^{-n}\mathcal O$ は $A$ 上平坦である。
AZ の生成条件により任意の coherent 対象はこれらの有限直和の商である。
従って任意の有限提示 $A$-加群 $M$ について
$\operatorname{Tor}^A_1(M,-)$ は co-effaceable となる。
これが \cite[Proposition 3.1, Definition 3.2]{LowenVanDenBergh2006}
の flatness 条件である。Ind 圏についても同論文 Proposition 3.6 を適用できる。

自然自己同型 $\eta$ の $s^{-n}\mathcal O$ 上の成分は、相対例外性により
単元 $a_n\in A^\times$ の scalar である。
正則二次切断の完全列により
\[
 \operatorname{Hom}(s^{-n}\mathcal O,\mathcal O)
       \simeq (\mathcal B/(g))_n
\]
は階数 $(n+1)^2$ の有限局所自由 $A$-加群、従って忠実である。
自然性は $(a_n-a_0)(\mathcal B/(g))_n=0$ を与えるので $a_n=a_0$。
AZ の全射と自然性から全ての対象上で $\eta=a_0\operatorname{id}$ となる。
pointed 二射では $a_0=1$ である。
この結論は fppf 局所で確認できるため、根を大域に選んでいない族にも成り立つ。
\end{proof}

